\chapter*{Abstract}
\addcontentsline{toc}{chapter}{Abstract}
Optical flow estimation in 2D anime is difficult due to large flat-color regions with almost no texture, sharp contour lines that create discontinuous motion boundaries, multi-layer occlusions, and stylized or exaggerated character motions that break the assumptions of natural-image flow models. Unsupervised two-frame approaches therefore drift in uniform areas and fail to maintain stable correspondences over time. We propose AniFlowFormer-T, an unsupervised multi-frame temporal transformer designed specifically for the motion structure of 2D anime. The model aggregates correspondence cues across a temporal window, maintains long-range stability through a latent cost memory, and refines motion via a global temporal regressor, producing flow estimates that remain reliable even in texture-sparse or rapidly changing sequences. Without any ground-truth labels, AniFlowFormer-T reduces the average EPE on AnimeRun from 4.21 (RAFT) to 3.52 and lowers large-motion error by 34\%, and when integrated into the ATD-12K interpolation pipeline, it improves PSNR from 31.01 dB to 31.54 dB. These results show that temporal reasoning and long-range aggregation are essential for accurate optical flow in 2D anime.
